\chapter{Installation, configuration and execution}\label{ch:installation}
\section{Prerequisites and source distribution}
Use an existing authorized checkout or source archive with its maintained
shared libraries, inputs and test fixtures. Preserve the supplied directory
layout. This edition does not prescribe an obsolete CVS host or silently
retrieve missing model data. Record the source revision and local changes
with each production build.

The native toolchain needs a compatible C++ compiler, make, Perl and MPI;
dependency-light modern contracts commonly require C++17. Python 3 runs the
global test procedure and many post-processors. A Fortran compiler and matching
MPI interface are required for SWMF and selected linked modules. Optional
capabilities can require CSPICE, external field libraries, numerical tables,
matplotlib/NumPy or an accelerator toolchain. Inspect the chosen application
and machine profile rather than assuming all dependencies are universal.

\section{Site configuration and application selection}
The established \path{Config.pl} workflow selects the application and site
settings. It uses \path{input/NAME.input} first and falls back to
\path{input/NAME.conf} only if the former does not exist. Initial installation
and compiler selection are site-specific; inspect \code{./Config.pl -help}
and the existing \path{Makefile.conf}. Preserve site libraries and settings.
Do not edit \path{.amps.conf} directly.

For a new compile-time application selection:
\begin{lstlisting}
./Config.pl -show
./Config.pl -application=enceladus
\end{lstlisting}
Changing the runtime .in file selects only among compiled capabilities; it
cannot change species IDs, dimension or allocated record layout.
Generated files in \path{build} or \path{srcTemp} are not the maintained
place to implement a lasting fix.

\section{Mandatory native rebuild sequence}
Before \emph{every} native AMPS rebuild, confirm the repository root and verify
that no build or test simulation is using it. In the current workspace the
root is \path{/home/vtenishe/Mars2/AMPS}. From that confirmed, idle root:
\begin{lstlisting}
# Remove only the generated root build directory.
rm -rf -- build

# Regenerate the chosen application's configuration/hooks.
./Config.pl -application=enceladus

# Use the established top-level compilation workflow.
env MAKEFLAGS="-j16" make amps
\end{lstlisting}
Some application guides explicitly stage generator hooks before make:
\begin{lstlisting}
./ampsConfig.pl -input enceladus.input -no-compile
\end{lstlisting}
Follow the selected adapter's established sequence. Do not remove source directories,
site settings, libraries or another checkout. A focused standalone compile of
a shared-model test is not a native AMPS build.

The global native test adapters perform and record this preflight under
exclusive ownership. Prefer them for reproducible verification. Do not run a
second manual Config.pl or make against their active checkout.

\section{SPICE configuration and data readiness}
SPICE-enabled applications require both the toolkit and appropriate kernels.
Configure explicit local paths through the established workflow, for example:
\begin{lstlisting}
./Config.pl -spice-path=/path/to/SPICE/cspice \
  -spice-kernels=/path/to/SPICE/Kernels
\end{lstlisting}
These paths are placeholders. A no-SPICE fixture selects
\code{-spice-path=nospice}; a component-level SPICE-off input by itself may
not remove an inherited global toolkit/link choice.
Test actual state and orientation queries for the intended target, frame and
epoch. The presence of a file with a plausible name is not proof of coverage.
Record versions/hashes and load order; never substitute an unrelated mission
ephemeris to make initialization succeed.

\section{Run directories, MPI and batch execution}
Use a fresh run directory and freeze the exact executable and inputs. Paths
inside application files may resolve relative to the input file or the run
directory; inspect the owner. Copy/include all dependent source distributions,
event files and data receipts, preserving that relationship.
A minimal launcher pattern, after preparing a complete run, is:
\begin{lstlisting}
cd /path/to/prepared/run
mpiexec -n 4 ./amps -input run.in > run.log 2>&1
\end{lstlisting}
This is a launch pattern, not a universal CLI or rank recommendation. I3 uses
\code{-i FILE -case N}; some legacy applications take no runtime file.
Choose the MPI rank flag and launcher from the site's configuration. Allocate
memory for blocks, ghosts, particles, node fields and output buffers; counts
of CPUs alone do not establish that a case fits.

A batch script should set the working directory, requested ranks/threads,
required modules and log destinations explicitly. It must retain the real
process return code. Do not assume a shell pipeline to a logging command
preserves the simulation's failure code. Use the scheduler allocation, not
an arbitrary host-specific job example from an old manual.

\section{A first verification before a science run}
Run the fast global campaign and then the relevant native/end-to-end owner:
\begin{lstlisting}
python3 -B src/test/run_tests.py \
  --config test/qualification-machine/passion.json \
  --campaign smoke --output-dir test_output/first-smoke
\end{lstlisting}
For another host, supply its real machine profile. Read the final individual
FAIL/ERROR/SKIP list and linked logs. Keep compiler/startup success separate
from completed physical stepping and scientific qualification.
\source{Config.pl; ampsConfig.pl; src/test/README.md; test/qualification-machine/README.md}
